\documentclass{llncs2e/llncs}
\usepackage{makeidx}  % allows for indexgeneration

\usepackage[pdftex]{graphicx}
\graphicspath{{images/}}

\usepackage{hyperref}
\usepackage{url} \urlstyle{sf}

\usepackage{color}
\def\todo#1{{\color{red}TODO\@: #1}}
\def\addref#1{{\color{red}$[$#1$]$}}

\usepackage{datetime}
\newdate{startdate}{15}{05}{2017}
\newdate{enddate}{04}{08}{2017}

\title{Challenges in Automated Test Data Generation for Dynamically Typed Programming Languages}
\titlerunning{Challenges in Automated Test Data Generation for Dynamically Typed Programming Languages}

\author{Simon Bihel\inst{1, 2}}
\authorrunning{Simon Bihel}
\institute{Universit\'e de Rennes 1, France \and \'Ecole Normale Sup\'erieure de Rennes, France\\
\email{\href{mailto:simon.bihel@ens-rennes.fr}{simon.bihel@ens-rennes.fr}}}

\begin{document}
\pagestyle{headings}  % switches on printing of running heads
\mainmatter%          % start of the contributions

\maketitle

\hrulefill%

\begin{center}
  \textbf{Supervisor:} Dr.\ Shin Yoo\\
  COINSE Lab, KAIST, South-Korea
\end{center}

\hrulefill%

\begin{center}
  \displaydate{startdate} -- \displaydate{enddate}
\end{center}

\vfill

\centerline{%
    \includegraphics[width=0.3\textwidth]{Logo_ENS_Rennes}
    \hspace{0.05\textwidth}
    \raisebox{1\height}{\includegraphics[width=0.3\textwidth]{Universite_Rennes_1_logo}}
    \hspace{0.05\textwidth}
    \includegraphics[width=0.3\textwidth]{KAIST_logo}
    \hspace{0.05\textwidth}
    \includegraphics[width=0.3\textwidth]{RB_NB}
}

\vfill

\begin{abstract}
Despite popular usage, most of the dynamically-typed languages still lack
automated test data generation tools. Most of the existing tools and approaches
depend critically on static and explicit types, which makes it hard to port any
of them over to dynamically typed languages. Some projects have tried to tackle
the challenges of dynamically-typed languages or complex data structure but
without much success and without generating momemtum in the research community.
After trying to develop a better tool than existing ones we were able to refine
what problems are in the way of automated test data generation. This paper
exposes them and tries to give some clues on how to tackle them.
\keywords{software engineering, testing, coverage, experimental, dynamic typing}
\end{abstract}
%
\newpage


\section{Introduction}
Testing is an essential part of software development. Especially due to the lack
of formal specifications, tests are the pillars to ensure well-behavior of
programs, be it lack of crash or coherent results. Various metrics exist to
evaluate how thoroughly tested a software is. One of these is the \textit{code
coverage} that describes the degree to which the source code of a program is
executed when a particular test suite runs.

Automating the generation of tests is essential to free the programmer from
writing them as is it a time-consuming and error-prone task. Evolutionary
methods like search-based software testing have been of interest for several
decades~\cite{miller1976automatic}. In particular, the automated test data
generation problem~\cite{korel1990automated} which consists of generating inputs
to explore different kinds of properties for a given program.

Past works on test generation automation have been focused on statically-typed
languages. Having information on the type of test inputs allows tools to focus
on the values of inputs and the depth of complex structures. Having no knowledge
on the inputs adds a new step in the process of testing as we have to find
fitting types for parameters. It also induces limitations to the kind of bugs a
tool can detect as well as a heavier cost due to tries with wrongly typed
inputs.

While developing an automated test generator for Python (presented in
Section~\ref{contribution}) we came to the conclusion that we wouldn't have come
further than existing works (explained in Section~\ref{relatedwork}). Thus we
decide to write down (Section~\ref{challenges}) the challenges that the
community will have to tackle in order to have a fully automated test generator
for dynamically-typed languages and more generally dynamic languages.


\section{Background}
\label{relatedwork}

While a lot of work has been put in automated testing for statically typed
languages, dynamically typed languages have received very little attention. Only
a few projects exist and each has a different approach, different levels of
automation, and uses different features that usually ship with dynamic
languages.

Multiple reasons can explain this void. Culturally, languages like Python are
not under the same level of industrial scrutiny for structural coverage,
compared to C or Java. But with rise of significant projects (ranging from
popular web applications to scientific tools) written with dynamic languages the
light has to be moved toward this problem.  Despite this interest the automation
of tests generation is still a daunting task.

Before presenting others automated testing tools in~\ref{related_research} we
will first give an overview of search-based software engineering in~\ref{st},~\ref{sbse}
and explain some techniques used by practitioners for dynamic languages testing
in~\ref{techniques}.

\subsection{Software testing}
\label{st}

Testing consists of a System Under Test (SUT) (e.g.\ a whole program, a
function, an object's method, etc.) to which we will apply perturbation by
sending inputs and then analyze the output.

There are a broad variety of test and metric classes. Tests focus on one
particular aspect of software to be finer grained, faster to run, to be more
comprehensible in the bugs they find, have more precise goals to achieve for
developer, etc.

In automated testing, where a tool generates tests without the help of the
programmer, we can identify some test classes. We will focus on dynamic testing,
as opposed to static testing (e.g.\ static analysis). We have two kinds of
testing, black-box and white-box. We will focus on white-box but we will present
both.

\paragraph{Black-box testing} consists of analysing only the inputs and
outputs of the SUT\@. It is used to test functionalities, integration in other
systems or even purely random testing is useful to discover failure
domains~\cite{ahmad2014new}. A failure domain can be a range of values for which
the SUT does not have a satisfying response (e.g.\ crashes).

Because we are focusing on the problem of automated testing without much
information of what the SUT should output we will not be focusing on this kind
of testing.

\paragraph{White-box testing}, on the other hand, works with more information on
the SUT (e.g.\ access to its source code or runtime monitoring) we can have more
insightful feedbacks. We can isolate different units (i.e.\ parts) of the
software, analyze the data and control flow, or test only modified parts (by
what is called regression).  With access to the innards of the SUT we can make
sure that every parts of it are tested. We use a metric called \textit{code
coverage} to evaluate the scope of tests.

%BRANCH COVERAGE
Multiple coverage criteria exist: Has every function been called? Has each
statement been executed? Has each boolean sub-expression been evaluated both to
true and false? Has each branch of each control structure been executed? That
last criterion is called branch coverage and it is the one we are focusing one.
It is part of standards~\cite{BSI98,RTCA92} for quality assurance and safety.

In our case of dynamically-typed languages, working on code coverage seems to be
a good idea to make sure we are exploring every parts and functionalities of a
function as they might be dependent on the type of parameters. This particular
problem of exploring every branch of a function is called test data generation.
One way to tackle this problem of finding certain values to pass tests is to use
search-based techniques.

\subsection{Search-based software engineering}
\label{sbse}

Search-based software engineering~\cite{harman2001search} is the application of
metaheuristic algorithms to problem related to software engineering. Such
algorithms comprise genetic algorithms, hill climbing, etc. The goal is iterate
and find an optimum by executing for real the system under test or optimization.

When we are trying to explore different branches, we are trying to satisfy some
boolean conditions. We will thus try to find values that can reach and pass
these control structure. With feedback to know how close we are to go through a
condition, we can iterate by tweaking parameters and re-executing the SUT\@.
This way of solving appeared in 1990~\cite{korel1990automated} and is called
\textit{search-based test data generation}. It has been extensively applied to
branch
coverage~\cite{mcminn2004search,lakhotia2007multi,mcminn2007iguana,Kim2017ts}.

One example is the Alternating Variable Method~\cite{mcminn2016avmf} that
consists in playing with one parameter until no positive change is seen, and
then focus on another parameter.

% it is justified to use in our case because we cannot deduce everything from
% static analysis or symbolic execution

\subsection{Testing techniques} % in dynamic languages}
\label{techniques}

Following is a list of testing methods that the reader can refer to to
understand the scope of related tools.

\paragraph{Instrumentation} is the insertion of additional code in the SUT
during compilation or linking. It is often needed for white-box testing. This
code can provide various information like counters to record which statement are
executed or compute distances as needed in branch coverage search.

\paragraph{Mock objects} simulate to behavior of real objects. They are used to
have more control on the execution, e.g.\ to remove some side effects to test
one particular property of a program or simply to provide instrumentation. It is
a versatile tool vastly
used~\cite{mackinnon2000endo,taneja2010moda,freeman2004mock,tillmann2006mock}.
In the most extreme use cases they can be used to test only function calls for
example, and methods only return either null or another mock object.

\paragraph{Object testing} can be done by creating an instance, modifying its
state and then measuring some metrics. To modify the state of the instance a
method call sequence is defined. After that, a final method call is done and
metrics like code coverage or method calls can be measured.

\subsection{Research works}
\label{related_research}

% no \citeauthor or \citet with llncs :(
\paragraph{} Ducasse et al.~\cite{ducasse2011challenges} have already exposed
challenges and milestones in random testing of dynamically-typed languages.
They are talking about random testing but with a white-box approach. They also
bundle object-oriented with dynamic types problems. The challenges they
highlight are instance generation and execution, error identification, and
result understanding. Our work exposes more challenges as well as more practical
details which are explained in the next section. They also propose milestones
that we will discuss in Section~\ref{futureworks}. They propose static analysis
and dynamic type collection for type inference, regression from existing
instances, hand-written contracts for automated instances generation, and
sandboxing to protect the environment. In the end their work is different than
ours because we are primarily focusing on basic dynamic types and test coverage.

\paragraph{} RuTeG~\cite{mairhofer2011search} (Ruby Test case Generator) is a
Ruby tool that targets the dynamic and object-oriented features of Ruby (not
just dynamic types). The do object testing and do code (line) coverage.

The center of the tool is the test case generator. For the inputs it uses
generators, basic types (numbers, strings, objects, arrays) are supported
natively but for the rest it is to the user to write problem-specific
generators. That means that this tool is not fully automated (but the authors
hint at the re-usability of a set of generators~\cite{feldt2002biomimetic}). The
test case generator uses a Genetic Algorithm for each step of the test unit
(constructor, method call sequence, method under test). In general, an
individual will be composed of type patterns and according data generators.

Before starting the search, the source code is analyzed statically to reduce the
search space. Method names and argument lists information will be collected as
well as methods called on arguments. It will identify methods impacting the
state of the instance and also identify control structures (this is possible
because they ignore run-time code addition). It also disqualify inappropriate
types. The authors claim that the tool performs better than fully random
generation.

% (this study: complex data types, but in the end it's the user that has to
% write such complex generator)

% hard to know clearly what they achieved with these student projects
\paragraph{} SpLATS~\cite{splats} (SpLATS Lazy Automated Test System) is another
tool for Ruby that was written during a student project. It has a lazy approach
and uses mock objects for that purpose. It will collect information at runtime
on which methods should be implemented and will output most of the time other
mock objects or else primitives.

Only methods of classes constructed by standard methods can be tested. It is
assumed that there are no side effects and no interactions with the environment.
No inheritance. No blocks or variable sized arguments.

It constructs all possible sets of arguments with only nil or mocks. When
implementing a method, it will return either nil or another mock object.

\paragraph{} SPLAT~\cite{splat} (Simple Python Lazy Automated Tester) is a tool
for Python with a lazy approach and is another student project. It generates
unit tests for instruction coverage.  It works with bytecode and uses what is
called meta-parameters in Python to make mock objects. It uses a unit tests
library to have automatically generated assertions on the return value or
exceptions raised. It also gather structure information from the bytecode.

The tool does not handle internal methods (those surrounded by `\texttt{\_\_}')
which means it cannot capture primitive operations (e.g.\ addition). It also
excludes string predicates but points to a study~\cite{alshraideh2006search}
addressing string predicates with a search-based approach. It does not seem to
generate recursive parameters like lists of lists.

\paragraph{} These works were published at least 5 years ago but without
follow-up. While writing our tool, we realised there was hardly anything new but
maybe a potential different subset of cases handled. We thus decided to write in
details all problems that a perfect tool should solve and we aim at building
precise foundations to build on to in future works. As each of these projects
have different problematics depending on the language they are tackling, we aim
to identify some generalization for dynamically-typed languages. We use the tool
that we wrote to illustrate with examples the challenges seen in the next
section.


\section{Challenges}
\label{challenges}

This section presents challenges of testing dynamically-typed languages based on
previous research~\cite{ducasse2011challenges} and our experience. This list
can be added to an already long list of challenges for non dynamic languages
testing~\cite{mcminn2011search}.

\subsection{Pure dynamic types problems}
\paragraph{Type errors} One big limitation for an automated code coverage tester
is the inability to identify crashes. As we are trying to find coherent types
for parameters, if a type error is raised, how can we know it is the parameters
that are the problem and not the code? It distances us from the fault-based
testing approach which is especially useful for object-oriented programming
systems~\cite{hayes1994testing}.

\paragraph{Relations between parameters} Parameters search can get incredibly
complex if we think even beyond complex data structures, with for example
functions or indexes. Also an object can be constructed by a function. Because
of that, if we knew which parameters were involved in an operation that raised a
type error, we couldn't be sure that changing just these parameters would
improve things.

\paragraph{Working but unfit types} A parameter of a certain type might not
generate type errors but make the search much more difficult. If we have for
example a floating-point number where an integer then simply approaching a round
value might be difficult.

% \paragraph{Number of parameters} You can sometimes pass an arbitrary number of
% parameters to a function. In the end it's just a list...

\paragraph{Difficulty to reduce domain of possible types} A parameter can have
different types for different branches. Because of that, it is not possible to
reduce the possible types of a parameter for all branches. But it is possible
for one branch.

\subsection{The bigger picture}
Because dynamically-typed languages usually integrate other paradigms, like
object-oriented, we need to consider the impact of dynamic types. Also, known
challenges in static languages become even more problematic.
% non-problems?
% \paragraph{Instance generation and execution} When dealing with complex data
% structures like objects you might not have easy access to an initialization
% method.
% no constructors in smalltalk -> close to what we do with just lists? methods?
% still falls in our handling of type errors?
\paragraph{Execution} Controlling the environment of the SUT is essential for
performance and security purposes. Two problems arise. The first one is for the
language environment, such as primitives, as it can be modified dynamically by
programs. The second problem is about identifying side effects.

\paragraph{Understanding results} In a case outside of simple code coverage, you
might want to analyze the result and judge whether it is right. If a function
ends up working for different parameters types than the ones it was designed for
we are facing the oracle problem~\cite{barr2015oracle}.

\paragraph{Infinite search spaces} Many statically-typed languages have bounded
types by the number of bits. On the other hand many dynamically-typed languages
have infinite range for integer for example, or seemingly with an infinite value
being implemented by default. Having information on whether we are approaching a
condition becomes essential then.

\paragraph{Dynamic data structures} Generating data structures of variable size
has been the subject of many works.~\cite{lakhotia2008handling} uses both a
search-based approach and a set of constraint solving rules. Tools for
object-oriented software testing
like~\cite{tonella2004evolutionary,ciupa2008artoo} use available instances to
generate new instances. Arcuri and Yao~\cite{arcuri2008search} focused only on
containers (e.g.\ lists, trees, etc.) which are data structures designed to
store any arbitrary type of data, using search-based techniques. Addressing this
problem is also important for variable sized arguments as they can be seen as
lists or dictionaries/hash tables.

\paragraph{Dynamic insertion of code} All tools that we have seen exclude from
their scope functions that execute code contained in a string (e.g.\
`\texttt{eval}' in Python). One argument given is that it is not a commonly used
feature.  Nonetheless in the case of branch coverage, as we rely on
instrumentation to actual detect branch forks, one problem arise. We cannot have
unique identifiers for instrumentation function calls as the code might be
inserted or executed multiple time.

\paragraph{Objects state} Manipulating the state of an object instance is
difficult but necessary for testing. We have seen that it is generally done by
having a list of method calls. With languages that don't have constructors like
Smalltalk is it even impossible to do pure random testing.


\section{Contribution}
\label{contribution}

Our tool started as a fork of CAVM~\cite{Kim2017ts} which is a search based test
data generation tool for C. It uses the Alternating Variable Method and can
generate complex data structures.

Our tool handles primitive types and lists (that can contain different types and
lists). It basically encloses the AVM search catching type and index errors.

\subsection{Philosophy}
\label{philosophy}

Ideally we wanted to have a tool the most universal possible, i.e.\ without
Python specific features. We focused on primitive type (e.g.\ integers, strings)
with lists to have a complex data structure. Lists can have elements of
different types, even other lists.

When an index error occurs, because we don't have a distance function to tell
use how close we are to the index asked, we are mocking list parameters. We
decorate the access method and bundle with the index error exception the index
asked. We did this to avoid growing wrong lists in the dark because of a lack of
feedback. It also mean that when a index error occurs we are going to grow the
list instead of reducing another parameter if it is used as the index.

\subsection{Technical difficulties}
\label{tech_difficulties}

In Python, exceptions (i.e.\ errors) carry very little information. If the
exception arises from an operation involving two objects, what will be available
is the code line number, the two types involved and the values of the objects in
frame (i.e.\ function call).

In addition to that, we cannot really use mocked objects to decorate all methods
and catch errors. For example with an operation, it will check on one object if
the according method works, and then check the other object.

Not all classes can be inherited which is another limitation for mocking as
inheritance allows to keep the type. Even though it is discouraged to write
conditions on types (e.g.\ `\texttt{isinstance}' function in Python) it is still
something that is used.




% \section{Evaluation}
% \label{evaluation}
% \todo{}

% try on libraries that don't use complex data structures, like math libraries
% and such?


\section{Future Works}
\label{futureworks}

As the purpose of this paper was to establish goals in the aim of automated
testing for dynamically-typed languages, we will now plan on what to do next.

\subsection{Pursuing the original goal}

Studying the time spent on trying invalid types would give some hints on whether
it is viable to extend the search for very complex structures such as objects.
If we were to use decorator with a catch for exceptions for all objects methods,
it would be interesting to compute the overhead. This overhead would depend on
whether the SUT uses exceptions a lot as the overhead of the catch will be
noticeable only when an exception occurs.

We could compare our approach to~\cite{poulding2017automated} which uses the
multiple-dispatch feature of Julia (methods have different implementation to
handle different types for parameters).

Overall we should focus at first on a simple tool that works for every
situation. The actual search for values is secondary, but still needed to be
sure we explore every needed type for arguments.

% Evaluate with libraries that take as input only primitive types, for example
% math libraries.
% In practice, to find real world code to test that does not use objects we could

\subsection{Beyond random testing}

This section aims at giving clues to go beyond random testing and use previous
work done on dynamic languages.

Data-flow and fine grained runtime monitoring. Data-flow would allow us to know
which parameters have an impact on other variables or even on other parameters.
This would be useful if we knew which variables were involved in an operation
resulting in an error. Which in turn would require wrap every elementary
operation. We would have to use a wrapper because in the case of a language like
Python we cannot override basic types.

A lot of work has been done on regression
testing~\cite{gligoric2011smutant,haupt2011type,steinert2010continuous,yoo2012regression}.
Of course this is quite the opposite of the goal of automated test generation.
Having a starting set of inputs that we know are valid greatly reduces our main
problem of finding appropriate types. But our work on automated testing can be
used as an extension and enhancement of existing, hand-written tests.
For big projects, there should be preexisting tests for functionalities testing.
Knowledge can also be extracted by watching the whole program with the graph of
function calls.

Some languages like Python have type annotations. We could ask the programmer to
precise with annotations when parameters should be custom types and assume other
parameters are primitives.

On-the-fly type analysis~\cite{chambers1991iterative} and
verification~\cite{chugh2012nested} for dynamically-typed object-oriented
languages are possible but don't make sense without valid inputs.

Genetic Algorithms are possible as~\cite{bottaci2010type} shows how to produce
mutants for dynamically-typed languages by generating mutants using types
information available at run-time.


\section{Conclusion}
\label{conclusion}

In this paper we presented challenges in dynamically-typed languages software
testing. We saw that challenges arise from both the paradigm and the
implementation. Detailing and isolating down problems will allow more coherent
future works.


\section*{Acknowledgment}
Thanks to Shin Yoo and the whole COINSE team for being so friendly and
welcoming. Thanks also to KAIST for hosting me.

This internship was supported by the R\'egion Bretagne with the \textit{Jeune
\`a l'International} grant as well as the ENS Rennes with the \textit{Aide
Transport pour la Mobilit\'e \`a l'International} grant.

\bibliographystyle{llncs2e/splncs03}
\bibliography{biblio}

\end{document}
